% Part: first-order-logic
% Chapter: syntax-and-semantics
% Section: terms-formulas

\documentclass[../../../include/open-logic-section]{subfiles}

\begin{document}

\olfileid[hi]{fol}{syn}{frm}

\olsection{पद तथा \printtoken{P}{formula} का निर्माण}

प्रथम-क्रम भाषा~$\Lang L$ दिए जाने पर हम~$\Lang L$ की मूल शब्दावली से बनी
अभिव्यक्तियाँ परिभाषित कर सकते हैं। इनमें विशेषतः \emph{पद} और
\emph{!!{formula}s} सम्मिलित हैं।

\begin{defn}[पद]
\ollabel{defn:terms}
\emph{पदों} का समुच्चय~$\Trm[L]$, भाषा~$\Lang L$ के लिए निम्न प्रकार
आगमनात्मक रूप से परिभाषित है:
\begin{enumerate}
\item प्रत्येक !!{variable} एक पद है।
\item भाषा~$\Lang L$ का प्रत्येक !!{constant} एक पद है।
\item यदि $f$ कोई $n$-स्थानी !!{function} है और $t_1$, \dots, $t_n$
  पद हैं, तो $\Atom{f}{t_1, \ldots, t_n}$ एक पद है।
\tagitem{limitClause}{
  इसके अतिरिक्त कुछ भी पद नहीं है।}{}
\end{enumerate}
जिस पद में कोई !!{variable}s न हों, उसे \emph{बंद पद} कहते हैं।
\end{defn}

\begin{explain}
भाषा और पदों के हमारे विनिर्देशन में !!{constant}s प्रतीकों की एक अलग
श्रेणी के रूप में आते हैं, किंतु इसके बजाय उन्हें शून्य-स्थानी
!!{function}s में सम्मिलित किया जा सकता था। तब पदों की परिभाषा के दूसरे
खंड की आवश्यकता न होती। केवल $\Atom{f}{t_1, \ldots, t_n}$ को स्वयं $f$
समझना होता, यदि $n =
0$ हो।
\end{explain}

\begin{defn}[सूत्र]
\ollabel{defn:formulas}
\emph{!!{formula}s} का समुच्चय~$\Frm[L]$, भाषा~$\Lang L$ के लिए निम्न प्रकार
आगमनात्मक रूप से परिभाषित है:
\begin{enumerate}
\tagitem{prvFalse}{$\lfalse$ एक परमाण्विक !!{formula} है।}{}

\tagitem{prvTrue}{$\ltrue$ एक परमाण्विक !!{formula} है।}{}

\item यदि $R$ कोई $n$-स्थानी !!{predicate} है जो~$\Lang L$ में है, और
  $t_1$, \dots, $t_n$,~$\Lang L$ के पद हैं, तो $\Atom{R}{t_1,\ldots, t_n}$ एक
  परमाण्विक !!{formula} है।

\item यदि $t_1$ और $t_2$,~$\Lang L$ के पद हैं, तो $\Atom{\eq}{t_1, t_2}$
  एक परमाण्विक !!{formula} है।

\tagitem{prvNot}{यदि $!A$ !!a{formula} है, तो $\lnot !A$ भी
  !!a{formula} है।}{}

\tagitem{prvAnd}{यदि $!A$ और $!B$ !!{formula}s हैं, तो $(!A \land
  !B)$ भी !!a{formula} है।}{}

\tagitem{prvOr}{यदि $!A$ और $!B$ !!{formula}s हैं, तो $(!A \lor !B)$
  भी !!a{formula} है।}{}

\tagitem{prvIf}{यदि $!A$ और $!B$ !!{formula}s हैं, तो $(!A \lif !B)$
  भी !!a{formula} है।}{}

\tagitem{prvIff}{यदि $!A$ और $!B$ !!{formula}s हैं, तो $(!A \liff !B)$
  भी !!a{formula} है।}{}

\tagitem{prvAll}{यदि $!A$ !!a{formula} है और $x$ !!a{variable} है,
  तो $\lforall[x][!A]$ भी !!a{formula} है।}{}

\tagitem{prvEx}{यदि $!A$ !!a{formula} है और $x$ !!a{variable} है,
  तो $\lexists[x][!A]$ भी !!a{formula} है।}{}

\tagitem{limitClause}{इसके अतिरिक्त कुछ भी !!a{formula} नहीं है।}{}
\end{enumerate}
\end{defn}

\begin{explain}
पदों और !!{formula}s के समुच्चयों की परिभाषाएँ \emph{आगमनात्मक परिभाषाएँ}
हैं। मूलतः हम !!{formula}s का समुच्चय अनंत चरणों में बनाते हैं। आरंभिक चरण
में हम सभी परमाण्विक सूत्रों को सूत्र घोषित करते हैं; यह परिभाषा के आरंभिक
मामलों, अर्थात् इनके लिए दिए मामलों, के अनुरूप है:
\iftag{prvTrue}{$\ltrue$, }{}%
\iftag{prvFalse}{$\lfalse$, }{}%
$\Atom{R}{t_1,\dots,t_n}$ और $\Atom{\eq}{t_1,t_2}$। अतः ``परमाण्विक
!!{formula}'' का अर्थ इस रूप का कोई भी !!{formula} है।

परिभाषा के अन्य मामले पहले से निर्मित !!{formula}s से नए !!{formula}s बनाने
के नियम देते हैं। दूसरे चरण में हम इन नियमों से परमाण्विक !!{formula}s से
!!{formula}s बना सकते हैं। तीसरे चरण में हम परमाण्विक सूत्रों और दूसरे चरण
में प्राप्त सूत्रों से नए सूत्र बनाते हैं, और यह क्रम आगे चलता है।
!!{formula} वही और केवल वही है जो अंततः ऐसे किसी चरण में बनाया जाता है।
\end{explain}

रूढ़ि के अनुसार हम $\eq$ को उसके तर्कों के बीच लिखते हैं और कोष्ठक छोड़ देते
हैं: $\eq[t_1][t_2]$, $\Atom{\eq}{t_1,t_2}$ का संक्षेप है। इसके अतिरिक्त,
$\lnot \Atom{\eq}{t_1,t_2}$ को $\eq/[t_1][t_2]$ से संक्षिप्त किया जाता है।
सूत्र $(!B \ast !C)$ को लिखते समय, जो $!B$, $!C$ से दो-स्थानी
संयोजक~$\ast$ द्वारा बना है, हम प्रायः सबसे बाहरी कोष्ठक-युग्म छोड़कर
केवल~$!B \ast !C$ लिखेंगे।

\begin{intro}
कुछ तर्कशास्त्र-ग्रंथ यह अपेक्षा करते हैं कि !!{variable}~$x$,~$!A$ में
अवश्य आए, तभी
\iftag{prvEx}{$\lexists[x][!A]$ }{}%
\iftag{notprvEx,notprvAll}{}{और }%
\iftag{prvAll}{$\lforall[x][!A]$ }{}%
को
\iftag{notprvEx,notprvAll}{!!a{formula}}{!!{formula}s} माना जाए। यह अपेक्षा
न रखने से कोई समस्या नहीं होती और काम सरल हो जाता है।
\end{intro}

\begin{tagblock}{defNot,defOr,defAnd,defIf,defIff,defTrue,defFalse,defEx,defAll}
\begin{defn}
परिभाषित संक्रियकों से बने सूत्रों को इस प्रकार समझना है:

\begin{tagenumerate}{defTrue,defFalse,defNot,defOr,defAnd,defIf,defIff,defEx,defAll}

\tagitem{defTrue}{$\ltrue$ इसका संक्षेप है:
  \iftag{prvFalse}{$\lnot\lfalse$}{$(!A \lor \lnot !A)$, जहाँ कोई
    नियत परमाण्विक !!{formula}~$!A$ लिया गया है।}.}{}

\tagitem{defFalse}{$\lfalse$ इसका संक्षेप है:
  \iftag{prvTrue}{$\lnot\ltrue$}{$(!A \land \lnot !A)$, जहाँ कोई
    नियत परमाण्विक !!{formula}~$!A$ लिया गया है।}.}{}

\tagitem{defNot}{$\lnot !A$, $!A \lif \lfalse$ का संक्षेप है।}{}

\tagitem{defOr}{$!A \lor !B$ इसका संक्षेप है:
  \iftag{prvAnd}{$\lnot(\lnot !A \land \lnot !B)$}{$\lnot !A \lif
    !B$।}.}{}

\tagitem{defAnd}{$!A \land !B$ इसका संक्षेप है:
  \iftag{prvOr}{$\lnot(\lnot !A \lor \lnot !B)$}{$\lnot (!A \lif
    \lnot !B)$।}.}{}

\tagitem{defIf}{$!A \lif !B$ इसका संक्षेप है:
  \iftag{prvOr}{$\lnot !A \lor !B)$}{$\lnot (!A \land \lnot !B)$}.}{}

\tagitem{defIff}{$!A \liff !B$, $(!A \lif !B) \land (!B
  \lif !A)$ का संक्षेप है।}{}

\tagitem{defAll}{$\lforall[x][!A]$, $\lnot\lexists[x][\lnot !A]$ का संक्षेप है।}{}

\tagitem{defEx}{$\lexists[x][!A]$, $\lnot\lforall[x][\lnot !A]$ का संक्षेप है।}{}
\end{tagenumerate}
\end{defn}
\end{tagblock}

किसी विशिष्ट अनुप्रयोग की भाषा में काम करते समय हम प्रायः दो-स्थानी
!!{predicate}s और !!{function}s को उनके संबंधित पदों के बीच लिखेंगे।
उदाहरणतः अंकगणित की भाषा में $t_1 < t_2$ और $(t_1 + t_2)$, तथा समुच्चय
सिद्धांत की भाषा में $t_1 \in t_2$। अंकगणित की भाषा में उत्तरवर्ती फलन को तो
रूढ़ि से उसके तर्क के \emph{बाद} लिखा जाता है:~$t'$। किंतु औपचारिक रूप से ये
क्रमशः $\Atom{\Obj A^2_0}{t_1, t_2}$, $\Obj f^2_0(t_1, t_2)$,
$\Atom{\Obj A^2_0}{t_1, t_2}$ और $f^1_0(t)$ के रूढ़िगत संक्षेप मात्र हैं।

\begin{defn}[वाक्यविन्यासगत समानता]
प्रतीक $\ident$ बताता है कि दो प्रतीक-शृंखलाएँ विन्यास में समान हैं;
अर्थात् $!A \ident !B$ तभी और केवल तभी, जब $!A$ और $!B$ समान लंबाई की
प्रतीक-शृंखलाएँ हों और प्रत्येक स्थान पर उनमें एक ही प्रतीक हो।
\end{defn}

प्रतीक $\ident$ के दोनों ओर संयोजन से प्राप्त शृंखलाएँ रखी जा सकती हैं।
उदाहरणतः, $!A \ident (!B \lor !C)$ का अर्थ है: प्रतीक-शृंखला~$!A$ वही
शृंखला है जो इसी क्रम में आरंभिक कोष्ठक, शृंखला $!B$, प्रतीक $\lor$,
शृंखला~$!C$, और अंतिम कोष्ठक को जोड़ने से प्राप्त होती है। यदि ऐसा है, तो
हम जानते हैं कि $!A$ का पहला प्रतीक आरंभिक कोष्ठक है, $!A$ में $!B$ एक
उपशृंखला के रूप में (दूसरे प्रतीक से आरंभ होकर) आता है, उसके बाद~$\lor$
आता है, आदि।

चूँकि पद और !!{formula}s मूल अवयवों से आगमनात्मक परिभाषाओं द्वारा बनाए जाते
हैं, उनके बारे में कथन सिद्ध करने के लिए हम निम्न आगमन-सिद्धांतों का उपयोग कर
सकते हैं।

\begin{lem}[\emph{पदों पर आगमन का सिद्धांत}]
    \ollabel{lem:trmind}
    मान लें कि $\Lang L$ कोई प्रथम-क्रम भाषा है। यदि किसी गुणधर्म~$P$ के लिए
    %
    \begin{enumerate}
        \item वह प्रत्येक !!{variable}~$v$ के लिए लागू होता है,
        %
        \item वह प्रत्येक !!{constant}~$a$ के लिए लागू होता है जो~$\Lang L$ में है, और
        %
        \item वह $f(t_1,\dotsc,t_n)$ के लिए लागू होता है, जब भी वह
        $t_1$,~\dots, $t_n$ के लिए लागू होता है और $f$ कोई $n$-स्थानी
            !!{function} है जो~$\Lang L$ में है,
    \end{enumerate}
    जहाँ $t_1$,~\dots, $t_n$ को~$\Lang{L}$ के पद माना गया है, तो $P$,
    ~ $\Trm[L]$ के प्रत्येक पद के लिए लागू होता है।
\end{lem}

\begin{prob}
    \olref[fol][syn][frm]{lem:trmind} सिद्ध कीजिए।
\end{prob}

\begin{lem}[\emph{!!{formula}s पर आगमन का सिद्धांत}]
    \ollabel{thm:frmind}
    मान लें कि $\Lang L$ कोई प्रथम-क्रम भाषा है। यदि कोई गुणधर्म~$P$ सभी
    परमाण्विक !!{formula}s के लिए लागू होता है और वह ऐसा है कि
    %
    \begin{enumerate}
        \tagitem{prvNot}{वह $\lnot !A$ के लिए लागू होता है, जब भी
            $!A$ के लिए लागू हो;}{}
        \tagitem{prvAnd}{वह $(!A \land !B)$ के लिए लागू होता है,
            जब भी $!A$ और~$!B$ के लिए लागू हो;}{}
        \tagitem{prvOr}{वह $(!A \lor !B)$ के लिए लागू होता है,
            जब भी $!A$ और~$!B$ के लिए लागू हो;}{}
        \tagitem{prvIf}{वह $(!A \lif !B)$ के लिए लागू होता है,
            जब भी $!A$ और~$!B$ के लिए लागू हो;}{}
        \tagitem{prvIff}{वह $(!A \liff !B)$ के लिए लागू होता है,
            जब भी $!A$ और~$!B$ के लिए लागू हो;}{}
        \tagitem{prvEx}{वह $\lexists[x][!A]$ के लिए लागू होता है,
            जब भी~$!A$ के लिए लागू हो;}{}
        \tagitem{prvAll}{वह $\lforall[x][!A]$ के लिए लागू होता है,
            जब भी~$!A$ के लिए लागू हो;}{}
  \end{enumerate}
  जहाँ $!A$ और $!B$ को~$\Lang{L}$ के !!{formula}s माना गया है, तो $P$,
  ~ $\Frm[L]$ के सभी सूत्रों के लिए लागू होता है।
\end{lem}

\end{document}
